\documentclass[aspectratio=169,10pt]{beamer}
% Shared Beamer setup for Fall 2026 course slides.
\mode<presentation>{
  \usetheme{Madrid}
  \usecolortheme{default}
}

\definecolor{CalPolyGreen}{HTML}{154734}
\definecolor{CalPolyGold}{HTML}{BD8B13}
\definecolor{DeckInk}{HTML}{1F2933}
\definecolor{DeckMuted}{HTML}{5F6B7A}

\setbeamercolor{structure}{fg=CalPolyGreen}
\setbeamercolor{palette primary}{bg=CalPolyGreen,fg=white}
\setbeamercolor{palette secondary}{bg=DeckInk,fg=white}
\setbeamercolor{palette tertiary}{bg=CalPolyGold,fg=DeckInk}
\setbeamercolor{frametitle}{bg=CalPolyGreen,fg=white}
\setbeamercolor{title}{fg=CalPolyGreen}
\setbeamercolor{normal text}{fg=DeckInk,bg=white}
\setbeamercolor{block title}{bg=CalPolyGreen,fg=white}
\setbeamercolor{block body}{bg=black!3,fg=DeckInk}

\setbeamertemplate{navigation symbols}{}
\setbeamertemplate{footline}[frame number]
\setbeamertemplate{itemize item}{\color{CalPolyGold}$\blacktriangleright$}
\setbeamertemplate{itemize subitem}{\color{CalPolyGreen}$\bullet$}

\newcommand{\CourseNumber}{Course}
\newcommand{\CourseTitle}{Course Title}
\newcommand{\LectureNumber}{0}
\newcommand{\LectureTitle}{Lecture Title}
\newcommand{\LectureDate}{Date}
\newcommand{\CourseUnit}{Unit}

\newcommand{\coursenumber}[1]{\renewcommand{\CourseNumber}{#1}}
\newcommand{\coursetitle}[1]{\renewcommand{\CourseTitle}{#1}}
\newcommand{\lecturenumber}[1]{\renewcommand{\LectureNumber}{#1}}
\newcommand{\lecturetitle}[1]{\renewcommand{\LectureTitle}{#1}}
\newcommand{\lecturedate}[1]{\renewcommand{\LectureDate}{#1}}
\newcommand{\courseunit}[1]{\renewcommand{\CourseUnit}{#1}}

\author{Kirk Duran}
\institute{California Polytechnic State University, San Luis Obispo}
\date{\LectureDate}

\newcommand{\makedecktitle}{
  \begin{frame}[plain]
    \vspace{0.25in}
    {\usebeamerfont{title}\color{CalPolyGreen}\LectureTitle\par}
    \vspace{0.18in}
    {\large \CourseNumber{}: \CourseTitle\par}
    \vspace{0.08in}
    {\large Lecture \LectureNumber{} -- \LectureDate\par}
    \vfill
    {\small Kirk Duran \textbar{} Cal Poly, San Luis Obispo\par}
  \end{frame}
}

\newcommand{\placeholdernote}{
  \begin{block}{Placeholder}
    Replace these bullets with the final lecture material, examples, and in-class activities.
  \end{block}
}

\AtBeginSection[]{
  \begin{frame}{Roadmap}
    \tableofcontents[currentsection]
  \end{frame}
}

\graphicspath{{../assets/lecture-02/}}
\setbeamersize{text margin left=0.42in,text margin right=0.42in}

\coursenumber{CS 1001}
\coursetitle{Introduction to Computing}
\lecturenumber{2}
\lecturetitle{How Computers Execute Programs}
\lecturedate{August 26, 2026}
\courseunit{1. Computing and Program Execution}

\begin{document}

\makedecktitle

\begin{frame}{Today}
  \begin{itemize}
    \item What a computer is and why computers are useful.
    \item How computing developed from counting tools to modern electronic systems.
    \item How a computer stores programs, executes instructions, and represents information.
    \item How machine language, assembly, and higher-level programming languages connect.
    \item What computer science studies and where computers appear in society.
    \item Why computation has a physical cost in hardware, energy, and infrastructure.
  \end{itemize}
\end{frame}

\sectionframe{Part 1: What Is a Computer?}

\begin{frame}{What Are Computers?}
  \begin{columns}[T,onlytextwidth]
    \begin{column}{0.47\textwidth}
      \begin{cpdefinition}
        A computer is a programmable electronic device that processes data to perform tasks.
      \end{cpdefinition}
      \vspace{0.08in}
      \begin{itemize}
        \item \textbf{Input:} information enters the system.
        \item \textbf{Processing:} the CPU follows instructions and transforms information.
        \item \textbf{Memory/storage:} information is kept for immediate or later use.
        \item \textbf{Output:} the system produces a result or physical action.
      \end{itemize}
    \end{column}
    \begin{column}{0.48\textwidth}
      \imageplaceholder{computer-components.jpg}
    \end{column}
  \end{columns}
\end{frame}

\begin{frame}[shrink=10]{Why Do We Use Computers?}
  \begin{columns}[T,onlytextwidth]
    \begin{column}{0.53\textwidth}
      \begin{itemize}
        \item \textbf{Speed:} computers can perform enormous numbers of simple operations quickly.
        \item \textbf{Automation:} repetitive procedures can run with little human intervention.
        \item \textbf{Accuracy and repeatability:} the same instructions can be followed consistently.
        \item \textbf{Scale:} computers let us analyze, store, and communicate more information than a person could manage manually.
      \end{itemize}
      \begin{exampleblockcp}
        Applications include scientific simulation, genome analysis, CAD, robotics, gaming, communication, and business systems.
      \end{exampleblockcp}
    \end{column}
    \begin{column}{0.41\textwidth}
      \imageplaceholder{clerical-computation.jpg}
      \vspace{0.05in}
      \begin{keyidea}
        Computers amplify human procedures; they do not remove the need to decide which procedures make sense.
      \end{keyidea}
    \end{column}
  \end{columns}
\end{frame}

\sectionframe{Part 2: A Brief History of Computing}

\begin{frame}{From Counting Tools to Algorithms}
  \begin{columns}[T,onlytextwidth]
    \begin{column}{0.48\textwidth}
      \begin{block}{Antiquity}
        The abacus supported arithmetic by giving people a physical representation of quantities and operations.
      \end{block}
      \imageplaceholder{abacus.jpg}
    \end{column}
    \begin{column}{0.48\textwidth}
      \begin{block}{9th Century}
        Persian mathematician Muhammad ibn Musa al-Khwarizmi wrote systematic procedures for calculation. His name is the historical root of the word \emph{algorithm}.
      \end{block}
      \imageplaceholder{al-khwarizmi.jpg}
    \end{column}
  \end{columns}
  \vspace{0.08in}
  \begin{keyidea}
    Computing has always involved both \textbf{representations} of information and \textbf{procedures} for manipulating it.
  \end{keyidea}
\end{frame}

\begin{frame}{Mechanical Calculation and Programmable Machines}
  \begin{columns}[T,onlytextwidth]
    \begin{column}{0.47\textwidth}
      \begin{block}{17th Century}
        Blaise Pascal and Gottfried Wilhelm Leibniz developed mechanical calculators for arithmetic operations.
      \end{block}
      \imageplaceholder{mechanical-calculator.jpg}
    \end{column}
    \begin{column}{0.48\textwidth}
      \begin{block}{1830s}
        Charles Babbage designed the Analytical Engine as a general-purpose mechanical computer. Ada Lovelace described an algorithm intended for the machine.
      \end{block}
      \imageplaceholder{babbage-engine.jpg}
    \end{column}
  \end{columns}
\end{frame}

\begin{frame}[shrink=35]{1936: Turing's Abstract Computer}
  \begin{columns}[T,onlytextwidth]
    \begin{column}{0.53\textwidth}
      \begin{itemize}
        \item Alan Turing introduced a mathematical model now called a \textbf{Turing machine}.
        \item The model showed how a small set of simple operations could express arbitrary algorithms.
        \item The important idea is not the physical machine; it is a model of what computation itself can do.
      \end{itemize}
      \begin{cpdefinition}
        Informally, a system is \textbf{Turing complete} when it can simulate any computation that a Turing machine can perform, given enough time and memory.
      \end{cpdefinition}
      \begin{exampleblockcp}
        This is about computational expressiveness, not whether a system is convenient, fast, or intended to be used as a computer.
      \end{exampleblockcp}
    \end{column}
    \begin{column}{0.41\textwidth}
      \imageplaceholder{turing-machine.jpg}
    \end{column}
  \end{columns}
\end{frame}

\begin{frame}{1940s--1950s: Electronic Computers and Languages}
  \begin{columns}[T,onlytextwidth]
    \begin{column}{0.47\textwidth}
      \imageplaceholder{eniac.jpg}
      \begin{block}{1940s}
        Electronic machines such as ENIAC could perform complex calculations much faster than manual calculation.
      \end{block}
    \end{column}
    \begin{column}{0.47\textwidth}
      \imageplaceholder{punch-card.png}
      \begin{block}{1950s}
        Languages such as FORTRAN and Lisp changed how humans expressed programs, moving programmers farther from raw machine instructions.
      \end{block}
    \end{column}
  \end{columns}
\end{frame}

\begin{frame}{1970s--1990s: Computers Become Personal and Connected}
  \begin{columns}[T,onlytextwidth]
    \begin{column}{0.47\textwidth}
      \begin{block}{Personal Computing}
        Microprocessors made smaller, less expensive computers practical. Systems such as the Apple I helped move computing toward individual users.
      \end{block}
      \imageplaceholder{early-personal-computer.jpg}
    \end{column}
    \begin{column}{0.47\textwidth}
      \begin{block}{The Internet and Web}
        Networked computers transformed communication and data exchange. Tim Berners-Lee's World Wide Web made linked information much easier to publish and access.
      \end{block}
      \imageplaceholder{internet-network-map.jpg}
    \end{column}
  \end{columns}
\end{frame}

\begin{frame}[shrink=10]{21st Century: Mobile, Cloud, and AI}
  \begin{columns}[T,onlytextwidth]
    \begin{column}{0.53\textwidth}
      \begin{itemize}
        \item Computing moved from one machine on one desk to networks of devices and data centers.
        \item Cloud platforms let organizations rent computation and storage at large scale.
        \item Smartphones put general-purpose computers, cameras, sensors, and networks in billions of pockets.
        \item AI systems increasingly use large-scale computing for perception, prediction, language, and control.
      \end{itemize}
      \begin{keyidea}
        Modern computing is often a \textbf{system of computers}, not a single computer.
      \end{keyidea}
    \end{column}
    \begin{column}{0.41\textwidth}
      \imageplaceholder{ai-era.jpg}
    \end{column}
  \end{columns}
\end{frame}

\sectionframe{Part 3: How Computers Execute and Represent Information}

\begin{frame}{The Stored-Program Computer}
  \begin{columns}[T,onlytextwidth]
    \begin{column}{0.47\textwidth}
      \begin{keyidea}
        A modern general-purpose computer stores both \textbf{data} and \textbf{program instructions} in memory.
      \end{keyidea}
      \begin{itemize}
        \item \textbf{CPU:} executes instructions.
        \item \textbf{Memory:} holds information being actively used.
        \item \textbf{Input/output:} connects the machine to users, devices, files, and networks.
        \item \textbf{Storage:} keeps information when programs are not actively using it.
      \end{itemize}
    \end{column}
    \begin{column}{0.48\textwidth}
      \imageplaceholder{stored-program-architecture.jpg}
    \end{column}
  \end{columns}
\end{frame}

\begin{frame}[shrink=20]{Inside the CPU: Control, Arithmetic, and Registers}
  \begin{columns}[T,onlytextwidth]
    \begin{column}{0.50\textwidth}
      \begin{itemize}
        \item The \textbf{control unit} coordinates instruction execution.
        \item The \textbf{arithmetic/logic unit (ALU)} performs arithmetic and logical operations.
        \item \textbf{Registers} are tiny, very fast storage locations inside the processor.
        \item Main memory is much larger than the registers but typically slower to access.
      \end{itemize}
      \begin{exampleblockcp}
        A useful mental model: the CPU is the worker, registers are the worker's hands, memory is the nearby workbench, and long-term storage is the filing cabinet.
      \end{exampleblockcp}
    \end{column}
    \begin{column}{0.44\textwidth}
      \imageplaceholder{computer-components.jpg}
    \end{column}
  \end{columns}
\end{frame}

\begin{frame}[shrink=10]{The Instruction Cycle: Fetch, Decode, Execute}
  \begin{columns}[T,onlytextwidth]
    \begin{column}{0.46\textwidth}
      \begin{enumerate}
        \item \textbf{Fetch:} retrieve the next instruction from memory.
        \item \textbf{Decode:} determine what the instruction means.
        \item \textbf{Execute:} perform the requested operation.
        \item Store intermediate or final results where they are needed.
        \item Repeat.
      \end{enumerate}
      \begin{activity}
        Imagine the instruction is ``add the values in A and B and save the result in C.'' What must be fetched, processed, and stored?
      \end{activity}
    \end{column}
    \begin{column}{0.49\textwidth}
      \imageplaceholder{instruction-cycle.jpg}
    \end{column}
  \end{columns}
\end{frame}

\begin{frame}[shrink=10]{Do Computers ``Think''?}
  \begin{columns}[T,onlytextwidth]
    \begin{column}{0.53\textwidth}
      \begin{itemize}
        \item Traditionally, a CPU does not decide what its instructions \emph{should} mean. It executes operations defined by its architecture and program.
        \item This precision is useful: the same operation can be repeated consistently.
        \item It is also dangerous: flawed instructions can be followed just as faithfully as correct ones.
      \end{itemize}
      \begin{keyidea}
        Computers are physical systems that implement rules. Human judgment enters through the problems, programs, data, and goals we give them.
      \end{keyidea}
    \end{column}
    \begin{column}{0.40\textwidth}
      \imageplaceholder{lukyanov-water-computer.jpg}
      \begin{center}
        \scriptsize Lukyanov's hydraulic integrator reminds us that computation is not inherently electronic.
      \end{center}
    \end{column}
  \end{columns}
\end{frame}


\begin{frame}{Computer Memory: Locations and Addresses}
  \begin{columns}[T,onlytextwidth]
    \begin{column}{0.48\textwidth}
      \begin{itemize}
        \item Think of memory as a large collection of storage locations.
        \item Each location has an \textbf{address} that identifies where data can be found.
        \item Addresses are numbers; programmers often see them written in hexadecimal.
        \item Programs can store a value at an address and later retrieve it.
      \end{itemize}
      \begin{exampleblockcp}
        The address identifies the location. The value is the information stored at that location.
      \end{exampleblockcp}
    \end{column}
    \begin{column}{0.46\textwidth}
      \imageplaceholder{memory-addresses.jpg}
    \end{column}
  \end{columns}
\end{frame}

\begin{frame}[shrink=10]{Bits and Bytes Are Physical}
  \begin{columns}[T,onlytextwidth]
    \begin{column}{0.51\textwidth}
      \begin{cpdefinition}
        A \textbf{bit} is one binary digit: 0 or 1. A \textbf{byte} is typically 8 bits.
      \end{cpdefinition}
      \begin{itemize}
        \item Modern computers represent bits using physical electronic states, commonly implemented with transistors and related circuitry.
        \item The exact technology can change; the logical abstraction of 0 and 1 remains useful.
        \item Bytes can encode numbers, characters, portions of images, instructions, and many other kinds of information.
      \end{itemize}
    \end{column}
    \begin{column}{0.43\textwidth}
      \imageplaceholder{memory-cells-bits.jpg}
    \end{column}
  \end{columns}
\end{frame}

\begin{frame}[shrink=10]{Memory Scale and Memory Technology}
  \begin{columns}[T,onlytextwidth]
    \begin{column}{0.47\textwidth}
      \begin{block}{Scale}
        \begin{itemize}
          \item 1 byte = 8 bits
          \item KB: thousands of bytes
          \item MB: millions of bytes
          \item GB: billions of bytes
          \item TB: trillions of bytes
        \end{itemize}
      \end{block}
      \begin{keyidea}
        Large datasets and models can require enormous memory even though each bit stores almost nothing.
      \end{keyidea}
    \end{column}
    \begin{column}{0.47\textwidth}
      \imageplaceholder{magnetic-core-memory.jpg}\hfill
      \imageplaceholder{ram-module.jpg}
      \begin{block}{Technology}
        Earlier magnetic-core systems were physically large and expensive. Semiconductor memory stores vastly more information in far less space and at much lower cost per byte.
      \end{block}
      \imageplaceholder{memory-warehouse.jpg}
    \end{column}
  \end{columns}
\end{frame}

\begin{frame}{Why Binary?}
  \begin{columns}[T,onlytextwidth]
    \begin{column}{0.50\textwidth}
      \begin{itemize}
        \item Binary is a base-2 numeral system containing only 0 and 1.
        \item Digital hardware can reliably distinguish between two ranges of physical states.
        \item Transistors act like controlled electronic switches and are combined to build logic gates, memory, and processors.
      \end{itemize}
      \begin{exampleblockcp}
        Binary is not ``the language of electricity.'' It is a convenient logical representation implemented by physical devices.
      \end{exampleblockcp}
    \end{column}
    \begin{column}{0.43\textwidth}
      \imageplaceholder{transistor-switch.jpg}
    \end{column}
  \end{columns}
\end{frame}

\begin{frame}{Reading Binary Numbers}
  \begin{columns}[T,onlytextwidth]
    \begin{column}{0.56\textwidth}
      \begin{block}{Place values}
        From right to left, binary places represent 1, 2, 4, 8, 16, 32, and so on.
      \end{block}
      \begin{exampleblockcp}
        \texttt{0101} = 0(8) + 1(4) + 0(2) + 1(1) = 5.
      \end{exampleblockcp}
      \begin{activity}
        What decimal number is \texttt{1010}? What binary representation would you use for decimal 7?
      \end{activity}
    \end{column}
    \begin{column}{0.37\textwidth}
      \imageplaceholder{machine-code-cpu.jpg}
    \end{column}
  \end{columns}
\end{frame}

\begin{frame}{Binary Addition Works Like Familiar Addition}
  \begin{columns}[T,onlytextwidth]
    \begin{column}{0.45\textwidth}
      \begin{itemize}
        \item Add one column at a time.
        \item \texttt{0 + 0 = 0}
        \item \texttt{0 + 1 = 1}
        \item \texttt{1 + 1 = 10}: write 0 and carry 1.
        \item Carries propagate just as they do in base 10.
      \end{itemize}
      \begin{exampleblockcp}
        \texttt{0101} (5) + \texttt{0011} (3) = \texttt{1000} (8).
      \end{exampleblockcp}
    \end{column}
    \begin{column}{0.49\textwidth}
      \imageplaceholder{binary-addition.jpg}
    \end{column}
  \end{columns}
\end{frame}

\sectionframe{Part 4: From Machine Language to Python}

\begin{frame}{Machine Language: Instructions the Processor Executes}
  \begin{columns}[T,onlytextwidth]
    \begin{column}{0.52\textwidth}
      \begin{itemize}
        \item A processor understands a particular set of machine instructions.
        \item Machine code represents those instructions as bit patterns.
        \item Instructions can move data, perform arithmetic, compare values, jump to other instructions, and interact with hardware.
        \item Raw machine code is difficult for humans to write and inspect directly.
      \end{itemize}
    \end{column}
    \begin{column}{0.41\textwidth}
      \imageplaceholder{machine-code-cpu.jpg}
    \end{column}
  \end{columns}
\end{frame}

\begin{frame}{Assembly Language Gives Machine Instructions Names}
  \begin{columns}[T,onlytextwidth]
    \begin{column}{0.50\textwidth}
      \begin{cpdefinition}
        Assembly language is a low-level programming language whose instructions closely correspond to a processor's machine instructions.
      \end{cpdefinition}
      \begin{itemize}
        \item Mnemonics such as \texttt{MOV} and \texttt{ADD} are easier for humans than raw binary.
        \item An \textbf{assembler} translates assembly instructions into machine code.
        \item Different processor architectures have different instruction sets.
      \end{itemize}
    \end{column}
    \begin{column}{0.44\textwidth}
      \imageplaceholder{assembly-to-machine-code.jpg}
    \end{column}
  \end{columns}
\end{frame}

\begin{frame}[shrink=25]{Higher-Level Languages Automate the Translation Work}
  \begin{columns}[T,onlytextwidth]
    \begin{column}{0.50\textwidth}
      \begin{itemize}
        \item Early programmers sometimes translated symbolic instructions into machine code manually.
        \item Assemblers automated one layer of translation.
        \item Compilers and interpreters let programmers express algorithms using increasingly human-friendly abstractions.
        \item Languages such as FORTRAN, C, Java, and Python let us focus more on the problem and less on processor details.
      \end{itemize}
      \begin{keyidea}
        Programming languages are layers of abstraction over the operations the machine ultimately performs.
      \end{keyidea}
    \end{column}
    \begin{column}{0.43\textwidth}
      \imageplaceholder{compiler-automation.jpg}
    \end{column}
  \end{columns}
\end{frame}

\begin{frame}[shrink=10]{Why Learn to Code When Tools Can Generate Code?}
  \begin{columns}[T,onlytextwidth]
    \begin{column}{0.51\textwidth}
      \begin{itemize}
        \item Automation repeatedly removes some low-level work and creates new higher-level work.
        \item Today's AI tools can write boilerplate, summarize information, and help debug programs.
        \item Someone still has to define the problem, judge whether the result is correct, test edge cases, and understand consequences.
      \end{itemize}
      \begin{keyidea}
        The valuable skill is not merely typing syntax. It is being able to describe computation precisely enough to inspect, test, and control it.
      \end{keyidea}
    \end{column}
    \begin{column}{0.42\textwidth}
      \imageplaceholder{high-level-code.jpg}
    \end{column}
  \end{columns}
\end{frame}

\sectionframe{Part 5: Computer Science and Computer Systems}

\begin{frame}[shrink=20]{What Is Computer Science?}
  \begin{columns}[T,onlytextwidth]
    \begin{column}{0.48\textwidth}
      \begin{cpdefinition}
        Computer science studies computation: algorithms, data, programming languages, systems, and the limits and applications of computation.
      \end{cpdefinition}
      \begin{itemize}
        \item \textbf{Theory:} algorithms and computational complexity.
        \item \textbf{Systems:} operating systems, databases, networks, architecture.
        \item \textbf{Software:} programming languages, software engineering, design.
        \item \textbf{Applied computing:} AI, machine learning, graphics, robotics, security, data science, HCI, and more.
      \end{itemize}
    \end{column}
    \begin{column}{0.46\textwidth}
      \imageplaceholder{cs-fields-taxonomy.jpg}
    \end{column}
  \end{columns}
\end{frame}

\begin{frame}{Computer Science Changes With Society}
  \begin{columns}[T,onlytextwidth]
    \begin{column}{0.47\textwidth}
      \imageplaceholder{cs-degree-history.jpg}
      \begin{block}{Participation changes}
        Enrollment and degree patterns have risen and fallen with technology cycles, labor markets, culture, and access to computing education.
      \end{block}
    \end{column}
    \begin{column}{0.47\textwidth}
      \imageplaceholder{tech-jobs-history.jpg}
      \begin{block}{Jobs change too}
        Computing occupations evolve as technology changes. Historical charts are snapshots, not guarantees about future demand.
      \end{block}
    \end{column}
  \end{columns}
\end{frame}

\begin{frame}[shrink=20]{Who Uses Computers?}
  \begin{columns}[T,onlytextwidth]
    \begin{column}{0.50\textwidth}
      \begin{itemize}
        \item \textbf{General users:} communication, entertainment, school, work, finance, navigation.
        \item \textbf{Scientists:} simulation and data analysis.
        \item \textbf{Engineers:} embedded systems, CAD, control, testing.
        \item \textbf{Businesses:} databases, logistics, CRM/ERP, e-commerce.
        \item \textbf{Accessibility:} screen readers, communication aids, prosthetic control.
      \end{itemize}
      \begin{keyidea}
        Because computing reaches nearly everyone, privacy, security, accessibility, bias, and responsible design are technical concerns too.
      \end{keyidea}
    \end{column}
    \begin{column}{0.43\textwidth}
      \imageplaceholder{smartphone-users-2016-2022.jpg}
      \begin{center}\scriptsize Historical smartphone-adoption snapshot from the original lecture.\end{center}
    \end{column}
  \end{columns}
\end{frame}

\begin{frame}[shrink=10]{Are All Computers the Same?}
  \begin{columns}[T,onlytextwidth]
    \begin{column}{0.47\textwidth}
      \begin{block}{General-purpose}
        PCs and smartphones are designed to run many different programs.
      \end{block}
      \begin{block}{Specialized}
        GPUs, embedded controllers, and supercomputers may be optimized for particular workloads or environments.
      \end{block}
      \begin{block}{Distributed}
        Clusters and cloud systems divide work across many networked machines for scale and fault tolerance.
      \end{block}
    \end{column}
    \begin{column}{0.47\textwidth}
      \imageplaceholder{personal-computer-timeline.jpg}
    \end{column}
  \end{columns}
\end{frame}

\begin{frame}[shrink=35]{Quantum Computers and Emerging Directions}
  \begin{columns}[T,onlytextwidth]
    \begin{column}{0.52\textwidth}
      \begin{itemize}
        \item \textbf{Quantum computing:} qubits and quantum effects enable different algorithms; potential speedups apply to particular problem structures rather than every task.
        \item \textbf{Specialized hardware:} GPUs and other accelerators trade generality for efficiency on important workloads.
        \item \textbf{New architectures:} research explores neuromorphic, photonic, and other post-Moore approaches.
        \item \textbf{Software and systems:} AI frameworks, cloud systems, distributed computing, and edge computing spread computation across many kinds of devices.
        \item \textbf{Efficiency:} power consumption and cooling increasingly constrain what systems can do.
      \end{itemize}
      \begin{exampleblockcp}
        Shor's algorithm is a famous example where a sufficiently capable quantum computer would change the complexity of integer factoring.
      \end{exampleblockcp}
    \end{column}
    \begin{column}{0.42\textwidth}
      \imageplaceholder{quantum-computer.jpg}
    \end{column}
  \end{columns}
\end{frame}

\sectionframe{Part 6: Computation Has a Physical Cost}

\begin{frame}{Activity: Three Ways to Get a Times Table}
  \begin{activity}
    Suppose you want the complete multiplication table from 1 through 12. Which approach do you expect to use the most electricity?
  \end{activity}
  \vspace{0.12in}
  \begin{enumerate}
    \item Write a small nested loop and compute all 144 products locally.
    \item Search the web and retrieve a page containing the table.
    \item Ask a large AI system to generate the program and return the results.
  \end{enumerate}
  \vspace{0.12in}
  \begin{exampleblockcp}
    Ignore your own time for the moment. Think about the computation, networking, storage, and infrastructure involved behind each option.
  \end{exampleblockcp}
\end{frame}

\begin{frame}[shrink=30]{The Energy Comparison Is About Scale, Not Exact Numbers}
  \begin{columns}[T,onlytextwidth]
    \begin{column}{0.50\textwidth}
      \begin{itemize}
        \item A few arithmetic operations on a local processor require very little energy.
        \item A network query activates networking equipment, servers, storage, cooling, and other infrastructure.
        \item A large AI inference may perform far more computation than the tiny arithmetic task itself.
        \item Exact energy use depends on hardware, software, location, batching, model size, utilization, and many other details.
      \end{itemize}
      \begin{keyidea}
        Use the simplest tool that solves the problem well. Abstraction is valuable, but abstraction is not physically free.
      \end{keyidea}
    \end{column}
    \begin{column}{0.42\textwidth}
      \imageplaceholder{cloud-data-center.jpg}
    \end{column}
  \end{columns}
\end{frame}

\begin{frame}{Computing Depends on Physical Infrastructure}
  \begin{columns}[T,onlytextwidth]
    \begin{column}{0.50\textwidth}
      \begin{itemize}
        \item Data centers require electricity, networking, cooling, land, and hardware.
        \item Growing computing demand can affect decisions about power generation and transmission.
        \item This is another reason efficiency matters: faster or smaller algorithms can reduce real resource use at scale.
      \end{itemize}
      \begin{keyidea}
        Software looks abstract, but every program ultimately runs on physical machines somewhere.
      \end{keyidea}
    \end{column}
    \begin{column}{0.42\textwidth}
      \imageplaceholder{data-center-energy-headline.jpg}
    \end{column}
  \end{columns}
\end{frame}

\begin{frame}{Energy Production Is Also a Computing Problem}
  \begin{columns}[T,onlytextwidth]
    \begin{column}{0.46\textwidth}
      \begin{itemize}
        \item Electrical demand changes throughout the day.
        \item Renewable generation also changes with conditions such as sunlight and wind.
        \item Grid operators use sensing, forecasting, optimization, and control software to balance production and demand.
        \item The ``duck curve'' illustrates how net demand can fall during solar-heavy hours and rise rapidly later in the day.
      \end{itemize}
    \end{column}
    \begin{column}{0.49\textwidth}
      \imageplaceholder{duck-curve.png}
    \end{column}
  \end{columns}
\end{frame}


\begin{frame}[shrink=10]{Where Python Fits}
  \begin{columns}[T,onlytextwidth]
    \begin{column}{0.48\textwidth}
      \begin{block}{What you write}
        \texttt{print("Hello, world!")}\\
        \texttt{x = 5 + 3}
      \end{block}
      \begin{block}{What the computer eventually needs}
        Operations represented in forms the processor and operating system can execute.
      \end{block}
      \begin{keyidea}
        Python lets us work at a high level while software underneath handles many details of memory, machine instructions, files, and hardware.
      \end{keyidea}
    \end{column}
    \begin{column}{0.46\textwidth}
      \begin{block}{Before next time}
        \begin{itemize}
          \item Make sure you can open Python or Google Colab.
          \item Create, save, and find a notebook or source file.
          \item Remember: Python program $\rightarrow$ translated operations $\rightarrow$ machine execution.
        \end{itemize}
      \end{block}
      \begin{activity}
        Bring one question about something inside a computer that still feels mysterious.
      \end{activity}
    \end{column}
  \end{columns}
\end{frame}

\begin{frame}{Next Time}
  \begin{center}
    {\Huge Python Introduction}
  \end{center}
\end{frame}

\end{document}
